%\documentstyle [11pt]{mybook}
%\input{latexmac}
%
%\setcounter{chapter}{1}
%\setcounter{section}{3}
%\setcounter{page}{56}
%\begin{document}
\section{Otomatisasi}

Berikut adalah rutin PASCAL untuk menjalankan \( k\rho_i+\rho_j \)
pada matriks teraugmentasi.
\begin{center}
\begin{verbatim}
  PROCEDURE Pivot(VAR LinSys : AugMat;
                  k : REAL;
                  i, j : INTEGER)
  VAR
   Col : INTEGER;
  BEGIN
   FOR Col:=1 TO NumVars+1 DO
     LinSys[j,Col]:=k*LinSys[i,Col]+LinSys[j,Col];
  END;
\end{verbatim}
\end{center}
Tentu saja, ini hanyalah satu bagian dari sebuah program utuh, tetapi
contoh ini menunjukkan bahwa reduksi Gauss sangat cocok untuk pemrograman komputer.

Namun, ada beberapa kendala.
Sebagai contoh, beberapa kendala muncul karena komputer menggunakan
hampiran berpresisi terbatas untuk bilangan riil.

Sistem-sistem berikut memberikan contoh sederhana.
\begin{center}
  \setlength{\unitlength}{4pt}      % two pairs of lines
  \begin{picture}(65,30)(-5,-18)
     \thinlines
     \multiput(0,0)(40,0){2}{\begin{picture}(20,0)(0,0)
         \put(0,0){\line(0,1){10} }          % y-axis
            \put(0,0){\line(0,-1){10} }
            \put(-0.5,3){\line(1,0){1} }
            \put(-0.5,6){\line(1,0){1} }
            \put(-0.5,9){\line(1,0){1} }
            \put(-0.5,-3){\line(1,0){1} }
            \put(-0.5,-6){\line(1,0){1} }
            \put(-0.5,-9){\line(1,0){1} }
           %\put(-2,10){\makebox(0,0)[b]{\small \(y\)} }
         \put(0,0){\line(1,0){10} }         % x-axis
            \put(0,0){\line(-1,0){10} }
            \put(3,-0.5){\line(0,1){1} }
            \put(6,-0.5){\line(0,1){1} }
            \put(9,-0.5){\line(0,1){1} }
            \put(-3,-0.5){\line(0,1){1} }
            \put(-6,-0.5){\line(0,1){1} }
            \put(-9,-0.5){\line(0,1){1} }
           %\put(10,2){\makebox(0,0)[b]{\small \(x\)} }
        \end{picture} }

     \put(0,-13){\makebox(0,0)[t]{\small \(\begin{linsys}{2}
                                   x &+ &2y &= &3  \\
                                  3x &- &2y &= &1
                                 \end{linsys}\) }  }
     \put(40,-13){\makebox(0,0)[t]{\small \(\begin{linsys}{2}
                                   x &+ &2y &= &3  \\
                     1.000\,000\,01x &+ &2y &= &3.000\,000\,01
                                 \end{linsys}\) }  }
     \thicklines
       \multiput(9,0)(40,0){2}{\line(-2,1){11}  }     % lines
         \multiput(9,0)(40,0){2}{\line(2,-1){2.5}  }

       \put(2.4,2.4){\( \bullet \)  }
       \put(0,-1.5){\line(2,3){8}  }
         \put(0,-1.5){\line(-2,-3){4}  }
       \put(42.4,2.4){\( \bullet \)}     % lines
       \put(49.5,0){\line(-2,1){12}  }     % lines
         \put(48.5,0){\line(2,-1){3}  }

  \end{picture}
\end{center}
(Kedua garis pada gambar kedua sulit dibedakan.)
Kedua sistem tersebut memiliki \( (1,1) \) sebagai solusi tunggalnya.

Pada sistem pertama, perubahan kecil pada bilangan-bilangannya hanya akan
menghasilkan perubahan kecil pada solusinya:
\begin{equation*}
 \begin{linsys}{2}
  x  &+  &2y &= &3                  \\
  3x &-  &2y &= &1.008
 \end{linsys}
\end{equation*}
menghasilkan solusi \( (1.002,0.999) \).
Secara geometris, sedikit mengubah salah satu garis tidak banyak mengubah
titik potongnya.

Hal tersebut tidak berlaku untuk sistem kedua.
Perubahan kecil pada koefisien-koefisiennya
\begin{equation*}
 \begin{linsys}{2}
   x              &+ &2y  &=  &3                  \\
  1.000\,000\,01x &+ &2y  &=  &3.000\,000\,03
 \end{linsys}
\end{equation*}
menghasilkan jawaban yang sama sekali berbeda: \( (3,0) \).

Solusi contoh kedua
berubah drastis, bergantung pada digit ke-\( 9 \).
Ini merupakan kabar buruk bagi mesin yang menggunakan \( 8 \) digit untuk merepresentasikan bilangan riil.
Singkatnya, sistem yang hampir singular mungkin sulit
ditangani secara komputasi.

Hal lain yang dapat menimbulkan masalah adalah perambatan galat.
Dalam sistem dengan banyak persamaan (misalnya, 100 atau lebih), galat
pembulatan kecil pada tahap awal prosedur dapat terakumulasi hingga mendominasi
solusi akhir.

Masalah-masalah ini, dan banyak masalah lain yang serupa,
berada di luar cakupan buku ini, tetapi ingatlah bahwa
meskipun Metode Gauss selalu berhasil secara teori, meskipun
sebuah program menerapkan metode tersebut dengan benar,
dan meskipun jawabannya tercetak pada kertas bergaris hijau, semua itu tidak
berarti bahwa jawaban tersebut benar.
Dalam praktik, selalu gunakan paket perangkat lunak
yang dikembangkan oleh para ahli untuk mengatasi berbagai masalah yang mungkin timbul.
%\end{document}
